% GitHub Stars Fetcher Example
% Compile with: lualatex example.tex

\documentclass{article}
\usepackage[utf8]{inputenc}

% Enable shell escape for curl commands
% Compile with: lualatex --shell-escape example.tex

% Load all GitHub metrics at the beginning
\directlua{dofile("packages/github/github-all.lua")}

\begin{document}

\title{Awesome LuaTeX}
\author{timerring}
\date{\today}
\maketitle

% Table of Contents
\tableofcontents
\newpage

\section{Introduction}

Welcome to Awesome LuaTeX! This document showcases a collection of useful LuaTeX functions and packages that enhance your LaTeX documents with dynamic content and powerful features.

LuaTeX combines the typesetting power of TeX with the flexibility of the Lua programming language, enabling you to create interactive and data-driven documents.

\section{Packages}

This section contains various LuaTeX packages and utilities for enhancing your documents.

\subsection{GitHub}

This subsection contains utilities for integrating GitHub data into your LaTeX documents.

\subsubsection{GitHub Repository Stars}

Get the star count for a specific GitHub repository.

\paragraph{Introduction}

The GitHub Repository Stars fetcher displays the number of stars a specific repository has received.

\paragraph{Load Lua Script}

\begin{verbatim}
\directlua{dofile("packages/github/github-repo-stars.lua")}
\end{verbatim}

\paragraph{How It Works}

This script uses shields.io API to fetch repository star counts and formats them with "+" suffix for abbreviated numbers (e.g., "240k+").

\paragraph{Usage Instructions}

Use the \verb|\getgithubrepostars{owner/repo}| command:

\begin{verbatim}
\getgithubrepostars{facebook/react}
\end{verbatim}

\textbf{Note:} Compilation requires the \texttt{--shell-escape} option.

\paragraph{Usage Examples}

\begin{itemize}
    \item \textbf{React:} \getgithubrepostars{facebook/react} stars
    \item \textbf{Linux Kernel:} \getgithubrepostars{torvalds/linux} stars
\end{itemize}

\subsubsection{GitHub User Total Stars}

Get the total number of stars across all repositories for a GitHub user.

\paragraph{Introduction}

The GitHub User Total Stars fetcher calculates and displays the cumulative stars received across all of a user's repositories.

\paragraph{Load Lua Script}

\begin{verbatim}
\directlua{dofile("packages/github/github-user-stars.lua")}
\end{verbatim}

\paragraph{How It Works}

This script uses GitHub README Stats API to aggregate total stars from all user repositories, including public and private contributions.

\paragraph{Usage Instructions}

Use the \verb|\getgithubuserstars{username}| command:

\begin{verbatim}
\getgithubuserstars{torvalds}
\end{verbatim}

\textbf{Note:} This fetches total stars from all repositories owned or contributed to by the user.

\paragraph{Usage Examples}

\begin{itemize}
    \item \textbf{Linus Torvalds:} \getgithubuserstars{torvalds} total stars
    \item \textbf{timerring:} \getgithubuserstars{timerring} total stars
\end{itemize}

\subsubsection{GitHub Forks Counter}

Get the number of repository forks.

\paragraph{Introduction}

The GitHub Forks Counter fetches the fork count for any public GitHub repository.

\paragraph{Load Lua Script}

Load the script in your document:

\begin{verbatim}
\directlua{dofile("packages/github/github-forks.lua")}
\end{verbatim}

\paragraph{How It Works}

This script uses shields.io API to fetch fork counts and displays them in a readable format (e.g., "50k").

\paragraph{Usage Instructions}

Use the \verb|\getgithubforks{owner/repo}| command:

\begin{verbatim}
\getgithubforks{facebook/react}
\end{verbatim}

\paragraph{Usage Examples}

\begin{itemize}
    \item \textbf{React:} \getgithubforks{facebook/react} forks
    \item \textbf{Vue:} \getgithubforks{vuejs/vue} forks
\end{itemize}

\subsubsection{GitHub Watchers Counter}

Get the number of repository watchers.

\paragraph{Introduction}

The GitHub Watchers Counter fetches the watcher count for any public GitHub repository.

\paragraph{Load Lua Script}

\begin{verbatim}
\directlua{dofile("packages/github/github-watchers.lua")}
\end{verbatim}

\paragraph{How It Works}

Fetches watcher data from shields.io API and formats it appropriately.

\paragraph{Usage Instructions}

\begin{verbatim}
\getgithubwatchers{owner/repo}
\end{verbatim}

\paragraph{Usage Examples}

\begin{itemize}
    \item \textbf{React:} \getgithubwatchers{facebook/react} watchers
\end{itemize}

\subsubsection{GitHub Contributors Counter}

Get the number of repository contributors.

\paragraph{Introduction}

The GitHub Contributors Counter shows how many people have contributed to a repository.

\paragraph{Load Lua Script}

\begin{verbatim}
\directlua{dofile("packages/github/github-contributors.lua")}
\end{verbatim}

\paragraph{How It Works}

Fetches contributor count from shields.io API.

\paragraph{Usage Instructions}

\begin{verbatim}
\getgithubcontributors{owner/repo}
\end{verbatim}

\paragraph{Usage Examples}

\begin{itemize}
    \item \textbf{React:} \getgithubcontributors{facebook/react} contributors
\end{itemize}

\subsubsection{GitHub Issues Counter}

Get the number of open issues.

\paragraph{Introduction}

The GitHub Issues Counter displays the current number of open issues in a repository.

\paragraph{Load Lua Script}

\begin{verbatim}
\directlua{dofile("packages/github/github-issues.lua")}
\end{verbatim}

\paragraph{How It Works}

Fetches open issue count from shields.io API and extracts the number from "XXX open" format.

\paragraph{Usage Instructions}

\begin{verbatim}
\getgithubissues{owner/repo}
\end{verbatim}

\paragraph{Usage Examples}

\begin{itemize}
    \item \textbf{React:} \getgithubissues{facebook/react} open issues
\end{itemize}

\subsubsection{GitHub License Detector}

Get the repository license type.

\paragraph{Introduction}

The GitHub License Detector identifies the license under which a repository is distributed.

\paragraph{Load Lua Script}

\begin{verbatim}
\directlua{dofile("packages/github/github-license.lua")}
\end{verbatim}

\paragraph{How It Works}

Fetches license information from shields.io API and displays the license name (e.g., "MIT", "Apache-2.0").

\paragraph{Usage Instructions}

\begin{verbatim}
\getgithublicense{owner/repo}
\end{verbatim}

\paragraph{Usage Examples}

\begin{itemize}
    \item \textbf{React:} \getgithublicense{facebook/react}
    \item \textbf{Linux:} \getgithublicense{torvalds/linux}
\end{itemize}

\subsubsection{GitHub Release Version}

Get the latest release version.

\paragraph{Introduction}

The GitHub Release Version fetcher displays the most recent release tag.

\paragraph{Load Lua Script}

\begin{verbatim}
\directlua{dofile("packages/github/github-release.lua")}
\end{verbatim}

\paragraph{How It Works}

Fetches the latest release version from shields.io API (e.g., "v19.2.0").

\paragraph{Usage Instructions}

\begin{verbatim}
\getgithubrelease{owner/repo}
\end{verbatim}

\paragraph{Usage Examples}

\begin{itemize}
    \item \textbf{React:} \getgithubrelease{facebook/react}
\end{itemize}

\subsubsection{GitHub Tag Version}

Get the latest tag.

\paragraph{Introduction}

The GitHub Tag Version fetcher displays the most recent tag.

\paragraph{Load Lua Script}

\begin{verbatim}
\directlua{dofile("packages/github/github-tag.lua")}
\end{verbatim}

\paragraph{How It Works}

Fetches the latest tag from shields.io API.

\paragraph{Usage Instructions}

\begin{verbatim}
\getgithubtag{owner/repo}
\end{verbatim}

\paragraph{Usage Examples}

\begin{itemize}
    \item \textbf{React:} \getgithubtag{facebook/react}
\end{itemize}

\subsubsection{GitHub Language Detector}

Get the primary programming language.

\paragraph{Introduction}

The GitHub Language Detector identifies the main programming language used in a repository.

\paragraph{Load Lua Script}

\begin{verbatim}
\directlua{dofile("packages/github/github-language.lua")}
\end{verbatim}

\paragraph{How It Works}

Fetches language data from shields.io API and automatically capitalizes common language names (JavaScript, TypeScript, Python, etc.).

\paragraph{Usage Instructions}

\begin{verbatim}
\getgithublanguage{owner/repo}
\end{verbatim}

\paragraph{Usage Examples}

\begin{itemize}
    \item \textbf{React:} \getgithublanguage{facebook/react}
    \item \textbf{Linux:} \getgithublanguage{torvalds/linux}
\end{itemize}

\subsubsection{GitHub Creation Date}

Get the repository creation date.

\paragraph{Introduction}

The GitHub Creation Date fetcher displays when a repository was created.

\paragraph{Load Lua Script}

\begin{verbatim}
\directlua{dofile("packages/github/github-created.lua")}
\end{verbatim}

\paragraph{How It Works}

Fetches creation date from GitHub API and formats it as YYYY-MM-DD.

\paragraph{Usage Instructions}

\begin{verbatim}
\getgithubcreated{owner/repo}
\end{verbatim}

\paragraph{Usage Examples}

\begin{itemize}
    \item \textbf{React:} Created on \getgithubcreated{facebook/react}
\end{itemize}

\subsubsection{Complete Example}

Here's a comprehensive overview using all GitHub metrics:

\paragraph{Repository Metrics}

\begin{itemize}
    \item \textbf{Repository:} facebook/react
    \item \textbf{Repository Stars:} \getgithubrepostars{facebook/react}
    \item \textbf{Forks:} \getgithubforks{facebook/react}
    \item \textbf{Watchers:} \getgithubwatchers{facebook/react}
    \item \textbf{Contributors:} \getgithubcontributors{facebook/react}
    \item \textbf{Open Issues:} \getgithubissues{facebook/react}
    \item \textbf{License:} \getgithublicense{facebook/react}
    \item \textbf{Language:} \getgithublanguage{facebook/react}
    \item \textbf{Latest Release:} \getgithubrelease{facebook/react}
    \item \textbf{Latest Tag:} \getgithubtag{facebook/react}
    \item \textbf{Created:} \getgithubcreated{facebook/react}
\end{itemize}

\paragraph{User Metrics}

\begin{itemize}
    \item \textbf{User:} torvalds
    \item \textbf{Total Stars:} \getgithubuserstars{torvalds}
\end{itemize}

\end{document}

